% !TEX program = pdflatex
\documentclass[11pt]{article}
\usepackage[margin=1in]{geometry}
\usepackage{amsmath, amssymb}
\usepackage{booktabs}
\usepackage{listings}
\usepackage{xcolor}
\usepackage{tcolorbox}
\usepackage{hyperref}
\usepackage{enumitem}

\hypersetup{colorlinks=true, linkcolor=blue!60!black, citecolor=blue!60!black}

\lstset{
    basicstyle=\ttfamily\footnotesize,
    breaklines=true,
    breakatwhitespace=false,
    columns=fullflexible,
    keepspaces=true,
    showstringspaces=false,
    frame=single,
    framerule=0.4pt,
    rulecolor=\color{gray!50},
    backgroundcolor=\color{gray!5},
    xleftmargin=0pt,
    xrightmargin=0pt,
}

\newtcolorbox{promptbox}[1][]{
    colback=blue!3, colframe=blue!50!black, fonttitle=\bfseries,
    boxrule=0.5pt, arc=2pt, left=6pt, right=6pt, top=4pt, bottom=4pt, #1
}

\title{TileBench LLM Code-Generation Pipeline\\ \large Iterative Refinement with Roofline-Aware Stopping}
\author{TileBench Project}
\date{}

\begin{document}
\maketitle

\section{Overview}

We adapt KernelBench-style iterative refinement to the TileBench framework for
automated generation of high-performance GPU kernels in two target DSLs
(Triton and NVIDIA cuTile). For every operator $\mathcal{O}$ in the TileBench
suite, an LLM is given (i) the operator's natural-language description,
(ii) the operator's configuration file describing the input-shape sweep, dtype
sweep, and metric expressions, and (iii) the PyTorch reference implementation
that defines bit-equivalent ground truth. The LLM is asked to emit two
implementation files (\texttt{impl\_triton.py} and \texttt{impl\_cutile.py})
that conform to the TileBench harness contract. Each emission is evaluated by
running compile, correctness, and Proton-based GPU-kernel timing through the
harness. The numerical feedback is fed back into the prompt for the next
iteration. The loop terminates when either (a) ten iterations elapse, or (b)
the geometric mean of $\textsc{roofline\_pct}$ across the top-3 largest input
configurations per dtype reaches $0.80$.

\section{Pipeline}

\subsection{Per-iteration loop}

For iteration $i \in \{0, 1, \ldots, 9\}$ of a fixed operator $\mathcal{O}$:

\begin{enumerate}[leftmargin=*, label=\arabic*.]
    \item \textbf{Prompt construction}. Build a user message from the
    framework guide, the operator description, the operator's
    \texttt{config.yaml}, and the PyTorch reference
    \texttt{impl\_torch.py}. On $i > 0$, also include the previous
    iteration's two source files verbatim, a markdown table of the
    iteration trajectory so far (per-iter \textsc{stop\_score},
    \textsc{report\_geo}, verify status, and total wall-clock time), and
    if the previous iteration regressed or broke verification relative
    to the best verify-clean iter so far, also include that
    best-so-far's two source files as a reference rebuild point. Persist
    the prompt to \texttt{iter\_$i$/prompt.md}.

    \item \textbf{LLM call}. Submit \{system, user\} to the LLM
    (GPT-5.5 via the OpenAI Responses API or Claude Opus 4.7 via the
    Anthropic Messages API) with streaming enabled and reasoning effort
    fixed at the chosen level (\texttt{xhigh} or \texttt{high}). Wall
    time, response length, and token usage
    (\textsc{input\_tokens}, \textsc{output\_tokens},
    \textsc{reasoning\_tokens}) are recorded. The raw response is
    persisted to \texttt{iter\_$i$/response.md} and token usage to
    \texttt{iter\_$i$/llm\_usage.json}. If the cached response file is
    already present from a prior run, the LLM call is skipped.

    \item \textbf{Code extraction}. Parse two fenced code blocks
    (\texttt{impl\_triton.py} and \texttt{impl\_cutile.py}) from the LLM
    response and write them next to the prompt.

    \item \textbf{Evaluation} (subprocess, GPU-isolated). Import the two
    generated impl modules. For every $(backend, dtype, case)$ triple in
    the operator's case grid:
    \begin{enumerate}[label=\alph*., leftmargin=*]
        \item Generate inputs through the operator's input generator.
        \item Run the PyTorch reference to produce a ground-truth output.
        \item Call \texttt{impl.run(*inputs, autotune=True)} inside a
        \texttt{signal.alarm(900\,s)} guard so any single configuration
        that hangs is bounded. The triggered autotune sweep is timed.
        \item Verify the kernel's output against the PyTorch reference
        within the per-dtype tolerances specified in
        \texttt{config.yaml}'s \texttt{verify:} block (or the harness
        defaults).
        \item \emph{Anti-cache verify}: regenerate a second independent
        input set, run \texttt{impl.run} again, verify against the
        \emph{regenerated} reference. This catches the failure mode in
        which the impl Python-caches outputs keyed on tensor identity
        or version: a cached impl would return the previous output and
        fail this second verify.
        \item Time the kernel with Proton (\texttt{triton.profiler})
        using 5 warm-up calls and 20 measured calls. \textbf{Each
        warm-up and measured call is fed a freshly regenerated input
        tuple} (constructed by 25 invocations of the operator's input
        generator before the timed loop begins), which defeats any
        identity-, version-, or storage-pointer-keyed output cache.
        Proton aggregates the GPU-kernel time inside the
        \texttt{proton.scope("launch")} block and the
        \texttt{device\_type=CUDA} nodes of its hatchet tree, so the
        reported latency excludes Python wrapper, CUDA driver, and
        autotune-lookup overhead. L2 is flushed (64\,MB write) before
        every iteration of both the warm-up and the timed loop.
        \item Compute the arithmetic intensity
        $\text{AI} = \text{flops\_expr} / \text{bytes\_expr}$ at this
        case's parameters and dtype, the attainable roofline
        $\min(P_{\text{compute}}(dtype), \beta \cdot \text{AI})$ with
        $\beta$ the measured peak HBM bandwidth and $P_{\text{compute}}$
        the datasheet dense Tensor-Core throughput for the dtype, and
        $\textsc{roofline\_pct} =
        (\text{flops\_expr} / t) / \min(P_{\text{compute}}, \beta \cdot
        \text{AI})$.
    \end{enumerate}

    \item \textbf{Score aggregation}. Two scores are produced:
    \begin{itemize}[leftmargin=*]
        \item $\textsc{report\_geo}$: geometric mean of
        $\min(\textsc{roofline\_pct}, 1)$ across \emph{all}
        $(backend, dtype, case)$ triples. Used for final benchmark
        reporting.
        \item $\textsc{stop\_score}$: geometric mean of
        $\min(\textsc{roofline\_pct}, 1)$ restricted to the top-3
        largest cases per $(backend, dtype)$. Small-$N$ cases are
        launch-overhead-bound and can never hit $0.80$ of the roofline,
        so excluding them is necessary for the iteration loop to ever
        terminate by performance.
    \end{itemize}

    \item \textbf{Feedback persistence}. Write the full per-combo
    metrics, both scores, compile/autotune/verify error lists, and the
    LLM usage and per-phase timing breakdown to
    \texttt{iter\_$i$/feedback.json}.

    \item \textbf{Stopping check}. The loop stops if (i) the iteration
    is fully verify-clean, autotune-clean, and compile-clean and
    $\textsc{stop\_score} \geq 0.80$, or (ii) $i = 9$ (the 10-iter cap).
\end{enumerate}

\subsection{Best-iter promotion}

After the loop ends, the verify-clean iteration with the largest
$\textsc{stop\_score}$ is copied into a \texttt{final/} directory as the
canonical implementation for this operator. If no iteration was
verify-clean, the harness emits a warning and the highest-scoring (but
buggy) iteration is promoted with a flag.

\subsection{Hardware peak values (B200)}

The roofline ceilings are loaded from
\texttt{data/peak\_performance/B200.json}. Compute ceilings are the
NVIDIA HGX B200 datasheet \emph{dense} Tensor-Core throughputs (sparsity
disabled, since LLM-generated kernels do not prepare 2:4 sparse layouts):
fp16/bf16 = 2250\,TFLOPS, fp8 = 4500\,TFLOPS, int8 = 4500\,TOPS, TF32 =
1100\,TFLOPS, fp64 = 37\,TFLOPS. fp32 uses the TF32 Tensor-Core dense
value (1100\,TFLOPS) because Triton's \texttt{tl.dot} defaults to TF32
input precision. The HBM3e peak bandwidth is the measured value
$\beta = 6{,}539.4\,\text{GB/s}$ (not the datasheet $7{,}700\,\text{GB/s}$
nominal).

\section{System Prompt}

The system message is constant across operators and iterations and is read
verbatim from \texttt{tools/llm\_codegen/system\_prompt.md}:

\begin{promptbox}
\begin{verbatim}
You are an expert GPU kernel engineer fluent in Triton and NVIDIA cuTile.
You write production-quality kernels that:
  - Compile cleanly on Blackwell (sm_100, B200)
  - Verify bit-equivalently against a PyTorch reference within the
    tolerances given
  - Use Tensor Cores (tl.dot / ct.mma) whenever the operator is
    matmul-shaped
  - Use tiled loops + autotune over a small (<=30 cfgs) search space

You MUST follow the TileBench framework conventions provided in the user
message EXACTLY. Deviations cause silent harness failures.
\end{verbatim}
\end{promptbox}

\section{User Prompt --- Operator-Agnostic Template}

The user message is rebuilt for each iteration from a small number of
operator-specific files (description, config, PyTorch reference, and
the previous iteration's source code if any) plus a constant framework
guide that encodes the harness contract. The template below is what the
harness produces for any operator $\mathcal{O}$; tokens of the form
\texttt{\{\{...\}\}} are substituted at prompt-build time.

\subsection{Iteration 0 (initial) template}

\begin{promptbox}
\begin{verbatim}
# Task: implement operator `{{op}}` for TileBench

Follow the framework conventions below. Then implement the operator
described.

---

# TileBench Framework Conventions

{{framework_guide.md contents -- see Section 5}}

---

# Operator description: `{{op}}`

{{benchmarks/problems/current/{{op}}_current.md contents}}

---

# `config.yaml` for `{{op}}`

```yaml
{{benchmarks/operators/{{op}}/config.yaml contents}}
```

Available case variables: see `case_grid` and `case_defaults` above.
Your kernels must work for all combinations.

---

# PyTorch reference: `impl_torch.py`

Your generated kernels must match this reference's outputs within the
tolerances in `verify:` (or per-dtype defaults if absent).

```python
{{benchmarks/operators/{{op}}/impl_torch.py contents}}
```

---

## Output format (STRICT)

Return EXACTLY two fenced code blocks, in this order, with these titles:

    ```python title="impl_triton.py"
    # full Python file content here
    ```

    ```python title="impl_cutile.py"
    # full Python file content here
    ```

No other code blocks. No prose between or after the two blocks beyond a
1-2 sentence summary of your approach. Do NOT include test code, do NOT
include PyTorch reference code (that file is provided by the framework).
\end{verbatim}
\end{promptbox}

\subsection{Iteration $N > 0$ (feedback) template}

\begin{promptbox}
\begin{verbatim}
# Task: improve operator `{{op}}` for TileBench (iteration {{N}})

Your previous iteration ({{N-1}}) did not yet meet the stopping criterion
(`stop_score` >= 80%, computed on the **largest 3 cases per dtype**).
Read the trajectory below carefully -- if your last iteration
regressed vs the best verify-clean iter so far, you should consider
going back to that approach as your starting point and trying a
different optimization. Then re-emit BOTH files.

## Iteration trajectory so far

| iter | stop_score | report_geo | verify | autotune | notes |
|---|---:|---:|---|---|---|
| 0 | {{ss_0%}}  | {{rg_0%}}  | {{verify_status_0}} | {{at_0}} |     |
| 1 | {{ss_1%}}  | {{rg_1%}}  | {{verify_status_1}} | {{at_1}} |     |
| ... |            |            |                      |          |     |
| N-1 | {{ss_{N-1}%}} | {{rg_{N-1}%}} | {{verify_status_{N-1}}} | {{at_{N-1}}} |     |

{{ regression-warning sentence if applicable }}
{{ best-so-far summary if any iter was verify-clean }}

## Feedback from iteration {{N-1}}

### (compile errors, autotune timeouts, verify failures -- only the
sections that fired)

### Performance (geo-mean roofline = {{stop_score%}}, target >= 80%)

Per-(backend, dtype, case) detail, sorted worst first (top 20 listed):
- triton / fp16  / {{params}} : {{pct}}% of roofline ({{bound_by}})
- cutile / int8  / {{params}} : {{pct}}% of roofline ({{bound_by}})
- ...

Action: focus on the worst-performing combinations. If
bandwidth-bound, optimize memory coalescing / tile layout / async
copies. If compute-bound, ensure Tensor Cores (tl.dot / ct.mma) and
high arithmetic intensity per load.

---

## Your previous `impl_triton.py` (iteration {{N-1}})

```python
{{previous impl_triton.py contents}}
```

## Your previous `impl_cutile.py` (iteration {{N-1}})

```python
{{previous impl_cutile.py contents}}
```

{{ optional block, only included if N-1 regressed vs best-so-far: }}
---

## Best verify-clean iteration so far (iter {{best_iter}}) -- `impl_triton.py`

This iter scored `stop_score={{best_score%}}`. Your previous iteration
regressed from this. Use this code as the starting point if your
iter-(N-1) attempt broke things.

```python
{{best-so-far impl_triton.py contents}}
```

## Best verify-clean iteration so far (iter {{best_iter}}) -- `impl_cutile.py`

```python
{{best-so-far impl_cutile.py contents}}
```

---

# Reminders: TileBench Framework Conventions

{{ identical to the framework guide in the initial prompt }}

---

# Operator description, config.yaml, PyTorch reference

{{ identical content blocks to the initial prompt -- repeated each iteration
   so the LLM never needs to scroll back }}

---

## Output format (STRICT)

{{ identical instruction to the initial prompt }}
\end{verbatim}
\end{promptbox}

\section{Framework Guide (injected into every user message)}

The framework guide is a single 280-line markdown document
(\texttt{tools/llm\_codegen/framework\_guide.md}) that specifies the
harness contract. Editing this file is the single point of control for
how the LLM is taught the TileBench conventions; no Python source needs
to change. The guide covers:

\begin{itemize}[leftmargin=*]
    \item \textbf{Files to produce}: exactly two files
    \texttt{impl\_triton.py} and \texttt{impl\_cutile.py};
    \texttt{impl\_torch.py} is copied from the human-written reference
    and must not be regenerated.
    \item \textbf{Required exports}: \texttt{run(*args, autotune: bool =
    False, **kwargs)} whose signature matches the PyTorch reference's
    \texttt{run()} exactly, and \texttt{get\_last\_config() -> dict |
    None} for surfacing the autotune-winner configuration.
    \item \textbf{Autotune conventions}. Triton: use the
    \texttt{triton.autotune(configs=[...], key=[...])} decorator with
    $\leq 30$ configurations; read the autotune winner via
    \texttt{kernel.best\_config}. cuTile: use the project's
    \texttt{CutileAutotuner} helper with a module-level
    \texttt{\_SEARCH\_SPACE} list and a module-level
    \texttt{\_last\_autotune\_config: dict} mutated only via
    \texttt{.clear()} and \texttt{.update()} (the \texttt{global}
    keyword is explicitly disallowed). No
    \texttt{\_DEFAULT\_CONFIG}-style fallback path needs to be
    provided; autotune is always required. The launch-execution-control
    knobs \texttt{num\_warps} and \texttt{num\_stages} are optional in
    the Triton config space --- they should be swept only when the
    kernel's structure can plausibly benefit from them (e.g.\
    software-pipelined matmul gains from a \texttt{num\_stages} sweep;
    a one-pass element-wise kernel rarely does). When swept, reasonable
    ranges are $\texttt{num\_warps} \in \{2, 4, 8\}$ and
    $\texttt{num\_stages} \in \{1, 2, 3, 4\}$.
    \item \textbf{Choosing meaningful tile / BLOCK sizes}. The guide
    enumerates lower bounds, upper bounds, and recommended sweep ranges
    for each operator shape, to prevent the LLM from wasting autotune
    budget on configurations that cannot win. Tiles below the lower
    bounds (e.g.\ \texttt{BLOCK\_SIZE=128} for a 20\,M-element fp32
    pointwise kernel) miss vectorised loads and memory coalescing
    windows, while tiles above the upper bounds exceed B200's
    register / shared-memory budgets. Concretely:
    \begin{itemize}[leftmargin=*, itemsep=2pt]
        \item Lower bounds: 1D pointwise / element-wise
        $\texttt{BLOCK\_SIZE} \geq 512$; row-reduction / softmax / norm
        $\geq 256$ along the reduction axis; stencil / 2D conv $\geq 256$
        along the spatial inner axis; matmul / attention tile
        $\texttt{BLOCK\_M}, \texttt{BLOCK\_N} \geq 64$ and
        $\texttt{BLOCK\_K} \geq 32$ for Tensor-Core utilisation.
        \item Upper bounds: $\texttt{BLOCK\_SIZE} \leq 8192$ for
        fp16/fp32 pointwise ($\leq 16384$ for int8); matmul tile area
        $\texttt{BLOCK\_M} \times \texttt{BLOCK\_N} \leq 256 \times 256$
        for fp16/fp8/int8 ($\leq 128 \times 256$ for fp32 due to
        register pressure); $\texttt{BLOCK\_K} \leq 128$ for
        fp16/fp8/int8 ($\leq 64$ for fp32).
        \item Recommended sweep ranges (a starting point that fits in
        the 30-config budget): pointwise
        $\texttt{BLOCK\_SIZE} \in \{512, 1024, 2048, 4096\}$;
        per-row reduction
        $\texttt{BLOCK\_N} \in \{256, 512, 1024, 2048\}$;
        matmul
        $\texttt{BLOCK\_M}, \texttt{BLOCK\_N} \in \{64, 128, 256\}$,
        $\texttt{BLOCK\_K} \in \{32, 64, 128\}$. cuTile uses
        \texttt{tile}, \texttt{tm}, \texttt{tn}, \texttt{tk}, and
        \texttt{occupancy} as the analogous knobs (with
        $\texttt{num\_warps} \times \texttt{occupancy} \approx 64$ on
        B200).
    \end{itemize}
    \item \textbf{Multi-dtype support}. A single \texttt{run()} must
    handle every dtype listed in the config's \texttt{case\_grid.dtype}
    field --- preferably via Triton/cuTile dtype polymorphism, or, if
    needed, via per-dtype branches inside the Python wrapper rather
    than two separate kernels.
    \item \textbf{Hardware constraints} (B200, sm\_100). Triton TMA
    descriptors require the last block-shape dimension to occupy
    $\geq 16$ bytes; cuTile tile dimensions must be powers of two,
    with \texttt{padding\_mode=ZERO} (or \texttt{NEG\_INF} for
    max-reduction kernels) handling non-power-of-two problem sizes.
    Tensor-Core MMA is invoked via \texttt{tl.dot} or \texttt{ct.mma},
    with Triton's default \texttt{input\_precision="tf32"} promoting
    fp32 inputs to the TF32 TC pipeline.
    \item \textbf{Forbidden: cross-call output caching}. The harness
    calls \texttt{run()} repeatedly with both the same input objects
    (during warm-up and the timed loop) \emph{and} with freshly
    regenerated input objects (during the anti-cache verify step). Any
    Python-level cache keyed on tensor identity, version, or
    storage-pointer that returns a stored output instead of executing
    the kernel will be caught by the harness either as a verify
    failure (because the cache returns a stale output for the
    regenerated inputs) or as an over-roofline outlier (because the
    cached path takes near-zero time). The guide spells out this
    contract explicitly with the rationale that any kernel-launch
    bypass is treated as a reward-hacking attempt.
    \item \textbf{Common pitfalls}: importing unlisted modules, signature
    mismatch, returning a list when the reference returns a tensor,
    forgetting masks for non-power-of-two trailing dimensions, mutating
    inputs in-place, and returning the wrong type from
    \texttt{get\_last\_config}.
    \item \textbf{TL;DR checklist} at the bottom (one-line items the LLM
    can self-check before returning code).
\end{itemize}

\section{Evaluator Implementation Notes}

\subsection{Subprocess isolation}

Each iteration's evaluation runs in a dedicated Python subprocess
(\texttt{tools/llm\_codegen/evaluator\_runner.py}). The parent
orchestrator (\texttt{evaluator.py}) launches it with an overall
two-hour wall-clock cap and propagates the iter directory, operator
name, output JSON path, and per-case 15-minute autotune cap through
environment variables. Subprocess isolation is necessary because LLM
output occasionally segfaults, leaks CUDA memory, or hangs inside an
ill-formed autotune sweep; killing the subprocess restores a clean
state for the next iteration.

\subsection{Per-case wall-clock guard}

Inside the runner, each $(backend, dtype, case)$ triple is wrapped in
a \texttt{signal.SIGALRM} handler scheduled to fire after 15 minutes.
Any single configuration that triggers a non-terminating autotune
benchmark or deadlocks raises \texttt{TimeoutError} and is recorded
as an autotune failure for the affected backend, so a single
pathological configuration cannot starve the rest of the sweep.

\subsection{Anti-cache verification}

Before timing, two verify calls are issued back-to-back with two
independent freshly-generated input tuples; both must pass with the
operator's per-dtype tolerances. The second call is what catches any
identity- or version-keyed Python cache the LLM may have introduced
(see the forbidden-pattern section of the framework guide).

\subsection{Proton timing with rotating inputs}

Timing is performed in a custom wrapper around
\texttt{triton.profiler}. We pre-generate 25 fresh input tuples (5 for
warm-up, 20 for the timed loop), invoke L2 flush + the kernel
\textbf{once} per Proton \texttt{launch} scope, and parse the
hatchet-tree output to extract the GPU-only kernel time aggregated
across the 20 timed iterations. Rotating the inputs is the timing-loop
analogue of the anti-cache verify: an output cache that keys on tensor
identity or storage pointer cannot return a cached value for any of
the 20 distinct tuples and therefore cannot inflate the measurement
beyond what a real kernel would achieve.

\subsection{Recorded metrics}

\texttt{feedback.json} for each iteration records:

\begin{itemize}[leftmargin=*]
    \item \textsc{op}, \textsc{iter\_dir}, case counts (attempted /
    succeeded);
    \item \textsc{compile\_errors} (per backend), \textsc{autotune\_errors},
    \textsc{verify\_failures} (list of \{backend, dtype, params,
    error message\});
    \item \textsc{roofline\_per\_combo}: for every successfully timed
    triple, the autotune winner config, latency in milliseconds,
    arithmetic intensity, measured bandwidth, measured TFLOPS, the
    dtype-specific compute ceiling and bandwidth ceiling, the resulting
    $\textsc{roofline\_pct}$, and a string \textsc{bound\_by} indicating
    whether the roofline at this AI is bandwidth- or compute-bound;
    \item \textsc{report\_geo\_all\_combos} and
    \textsc{stop\_score\_top\_k\_per\_dtype} as defined above;
    \item \textsc{timing\_breakdown}:
    \texttt{evaluator\_total\_s},
    \texttt{autotune\_total\_s},
    \texttt{iter\_total\_s} (the wall-clock the orchestrator spent on
    this iteration including the LLM call);
    \item \textsc{llm\_usage}:
    \texttt{elapsed\_s},
    \texttt{input\_tokens},
    \texttt{output\_tokens},
    \texttt{reasoning\_tokens} (only for reasoning models), and
    \texttt{response\_chars}, extracted from the SDK response object
    after streaming completes.
\end{itemize}

\section{Run Summary}

After all iterations finish (or the stopping criterion fires), the
orchestrator writes \texttt{run\_summary.json} at the operator's
$(model, effort)$ directory containing the full per-iter history,
the best verify-clean iter, the best-any iter, and the promoted iter.
This is the single artifact the analysis stage consumes to produce
per-operator and aggregate plots.

\section{Output Directory Layout}

\begin{lstlisting}
benchmarks/llm_generated/{op}/{model}/{effort}/
  iter_0/
    prompt.md         # user message sent to LLM
    response.md       # full raw LLM response
    impl_triton.py    # parsed from response
    impl_cutile.py    # parsed from response
    impl_torch.py     # copied from benchmarks/operators/{op}/
    config.yaml       # copied from benchmarks/operators/{op}/
    llm_usage.json    # tokens and elapsed
    feedback.json     # all evaluator metrics
  iter_1/ ...
  ...
  iter_9/ ...
  final/              # copy of the promoted iter
  run_summary.json    # trajectory + best + promoted
\end{lstlisting}

\section{Cost and Time Considerations}

The \texttt{xhigh} reasoning effort level dedicates the bulk of each
LLM call to internal chain-of-thought reasoning rather than user-facing
output. Per-iteration wall times scale with prompt length (which grows
linearly with iteration index as the trajectory section and previous
code blocks accumulate) and with the operator's autotune-search-space
size (Triton and cuTile both compile every candidate configuration the
first time it is seen). For pointwise operators with small search
spaces the per-iteration cost is dominated by the LLM call (1-3
minutes at \texttt{high}, 5-15 minutes at \texttt{xhigh}); for
matmul-shaped operators with larger spaces the autotune sweep
dominates and can approach the 15-minute per-case cap. The full
45-operator sweep at \texttt{high} effort with 10 iteration cap is
estimated at $\sim$30 hours on a single B200 in our pilot runs;
\texttt{xhigh} approximately doubles this. Token-usage and per-phase
timing are recorded for every iteration so that a token-aware cost
metric (e.g.\ effective TFLOPS-per-thousand-tokens) can be reported
post-hoc.

\end{document}
